\label{chap:cobertura-taxonomia-numeros}
\index{número HMT!taxonomía}
\index{transductor Hensel!desplazamiento completo}
\index{número racional!periodicidad}

El \cref{chap:numero-medicion} definió un número HMT como una sección
compatible del sistema inverso de supervivencia, y no como una sucesión de
cifras desprovista de genealogía. Quedaban, sin embargo, dos preguntas que
conviene separar con precisión. La primera es una pregunta de
\emph{representación}: ¿qué valores reales puede representar el transductor
3-ádico antes de imponer los filtros particulares de cierre? La segunda es
una pregunta de \emph{selección}: ¿qué secciones del espacio bruto quedan
escogidas por APP, TRIT y TPK cuando se exigen simultáneamente sus
compatibilidades? Esta sección resuelve la primera, formula la segunda como
una restricción explícita de dominio y extrae de la carta de base mil una
taxonomía completa de los valores visibles. Las herramientas de base son la
aritmética \(p\)-ádica, la dinámica simbólica y la caracterización posicional
de los racionales \cite{Koblitz1984,LindMarcus1995,HardyWright2008}.

La distinción es estructural. Un lenguaje capaz de representar toda cinta no
selecciona por ese solo hecho la cinta de \(\pi\), de \(e\), de \(\varphi\) o
de \(\alpha\). A la inversa, una ley de selección no puede evaluarse si antes
no se ha construido el espacio en el que actúa. Representación y selección
son dos teoremas consecutivos, no dos nombres para un mismo paso.

\section{Transductor de Hensel no filtrado y desplazamiento completo sobre las fibras ternarias}

Sea \(p\) un primo, sea \(d\geq1\) y sea
\(A\in\operatorname{GL}_d(\mathbb Z)\). Para un vector fila
\(x_t\in\mathbb Z_p^d\), definimos
\begin{equation}
 y_t=x_tA,
 \qquad
 u_t=y_t\bmod p\in\mathbb F_p^d,
 \qquad
 x_{t+1}=\frac{y_t-u_t}{p}.
 \label{eq:transductor-hensel-general}
\end{equation}
El bloque \(u_t\) es la emisión visible y \(x_{t+1}\) es la memoria que
permanece después de retirar esa cifra.

\begin{theorem}[Conjugación Hensel con el desplazamiento completo]
\label{thm:shift-hensel-completo}
Para todo \(T\geq1\), la aplicación
\[
 \Psi_T:(\mathbb Z/p^T\mathbb Z)^d
 \longrightarrow(\mathbb F_p^d)^T,
 \qquad
 x_0\longmapsto(u_0,\ldots,u_{T-1}),
\]
es biyectiva. Su inversa es
\begin{equation}
 x_0\equiv
 \sum_{t=0}^{T-1}p^t u_tA^{-(t+1)}
 \pmod {p^T}.
 \label{eq:inversa-hensel-finita}
\end{equation}
Al pasar al límite inverso se obtiene un homeomorfismo
\[
 \Psi:\mathbb Z_p^d\xrightarrow{\ \sim\ }
 (\mathbb F_p^d)^{\mathbb N}.
\]
Si \(S_A\) denota la actualización \(x_t\mapsto x_{t+1}\) y
\(\sigma\) elimina el primer bloque de una cinta, entonces
\[
 \Psi\circ S_A=\sigma\circ\Psi.
\]
\end{theorem}

\begin{proof}
De \cref{eq:transductor-hensel-general} se obtiene
\[
 x_t=(u_t+px_{t+1})A^{-1}.
\]
La iteración de esta identidad da
\[
 x_0=
 \sum_{t=0}^{T-1}p^t u_tA^{-(t+1)}
 +p^Tx_TA^{-T}.
\]
El último sumando se anula módulo \(p^T\), lo que prueba que una palabra de
\(T\) bloques determina a lo sumo una clase inicial. Recíprocamente, dada
cualquier palabra, se fija \(x_T=0\) y se reconstruye hacia atrás mediante
\(x_t=(u_t+px_{t+1})A^{-1}\). La igualdad
\(x_tA=u_t+px_{t+1}\) verifica que la clase obtenida emite la palabra
prescrita. Por tanto \(\Psi_T\) es biyectiva.

Las biyecciones son compatibles con la reducción módulo \(p^T\). Sus límites
inversos producen el homeomorfismo \(\Psi\); equivalentemente, la serie de
\cref{eq:inversa-hensel-finita} converge en \(\mathbb Z_p^d\). Aplicar
\(S_A\) retira exactamente el bloque \(u_0\), de modo que corresponde a
aplicar \(\sigma\) en la cinta.
\end{proof}

En la instancia HMT, \(p=3\), \(d=6\) y la matriz publicada tiene
determinante uno. Por ello
\begin{equation}
 \mathbb Z_3^6\simeq(\mathbb F_3^6)^{\mathbb N},
 \qquad |\mathbb F_3^6|=3^6=729.
 \label{eq:hmt-shift-729}
\end{equation}
Cada cinta de bloques de seis trits posee una única memoria 3-ádica inicial.
La matriz determina la coordenatización y la dinámica de lectura; la
universalidad del alfabeto demuestra que no excluye ninguna cinta.

\begin{corollary}[Entropía y medida de Haar]
La actualización Hensel bruta es conjugada al desplazamiento unilateral
completo de
\(729\) símbolos. Su entropía topológica es \(6\log3\). La medida de Haar de
\(\mathbb Z_3^6\) se transporta a la medida producto uniforme sobre los
bloques.
\end{corollary}

\begin{proof}
Existen \(729^T=3^{6T}\) cilindros de longitud \(T\), y cada clase módulo
\(3^T\) tiene masa de Haar \(3^{-6T}\). La conjugación del
\cref{thm:shift-hensel-completo} transporta ambas propiedades al
desplazamiento.
\end{proof}

La conclusión probabilística se refiere al espacio bruto dotado de Haar. No
afirma que las secciones publicadas de las constantes hayan sido muestreadas
según esa medida ni que sean algorítmicamente aleatorias.

\section{Carta ternaria balanceada de los enteros}

La carta 3-ádica bruta y la carta arquimediana no agotan los tipos numéricos
de la clasificación HMT. Para el operador entero se utiliza la palabra vacía junto con las
palabras reducidas
\[
 \mathscr W_{\mathrm{bal}}
 :=\{\varnothing\}\cup
 \{a_1\cdots a_m:
 a_j\in\{-1,0,+1\},\ a_1\ne0\}.
\]
Definimos recursivamente
\[
 b(\varnothing)=0,
 \qquad
 b(wa)=3b(w)+a.
\]

\begin{theorem}[Representación ternaria balanceada única]
\label{pf84:C03}
La evaluación
\[
 b:\mathscr W_{\mathrm{bal}}\longrightarrow\ZZ
\]
es biyectiva. Las palabras cuyo primer trit es \(+1\) representan los
enteros positivos y las que comienzan por \(-1\), los negativos.
\end{theorem}

\begin{proof}
Para \(n\ne0\), el último trit \(a\) es el representante único de
\(n\bmod3\) en \(\{-1,0,+1\}\). El padre
\[
 \frac{n-a}{3}
\]
tiene valor absoluto estrictamente menor. La iteración termina en \(0\), de
modo que todo entero posee una expansión finita. La unicidad del residuo
balanceado en cada paso prueba la unicidad de la palabra. El signo queda
fijado por el primer trit no nulo.
\end{proof}

\begin{remark}
Esta biyección es una carta aritmética exacta. No identifica todas las
palabras balanceadas con secciones supervivientes TPK ni sustituye la
genealogía, el acarreo o la orientación de una ruta enriquecida.
\end{remark}

\section{Sobreyectividad de la carta arquimediana no filtrada sobre \(\mathbb R\)}

Concatenemos las coordenadas de los bloques
\(u_t\in\mathbb F_3^6\) para obtener una cinta
\((a_n)_{n\geq1}\in\{0,1,2\}^{\mathbb N}\). Su evaluación ternaria es
\[
 \operatorname{ev}_3((a_n))=
 \sum_{n\geq1}\frac{a_n}{3^n}\in[0,1].
\]
Adoptamos la convención canónica que excluye la representación con una cola
finalmente constante de doses sólo cuando el mismo punto posee otra expansión
en la carta. El extremo \(1\) conserva, por tanto, su única representación
fraccionaria \(0.222\ldots{}_3\).

\begin{theorem}[Cobertura arquimediana bruta]
\label{thm:cobertura-arquimediana-bruta}
La aplicación
\[
 P_{\rm raw}:\mathbb Z_3^6
 \xrightarrow{\Psi}(\mathbb F_3^6)^{\mathbb N}
 \xrightarrow{\operatorname{flat}}\{0,1,2\}^{\mathbb N}
 \xrightarrow{\operatorname{ev}_3}[0,1]
\]
es continua y sobreyectiva. Por tanto,
\[
 \mathbb Z_3^6/\!\sim_{P_{\rm raw}}
 \ \cong\ [0,1]
\]
como espacio de valores, donde \(x\sim_{P_{\rm raw}}y\) si sus evaluaciones
coinciden.
\end{theorem}

\begin{proof}
Toda cifra ternaria pertenece a un único bloque consecutivo de seis cifras,
de modo que \(\operatorname{flat}\) es una biyección de espacios de cintas.
Toda \(r\in[0,1]\) posee una expansión ternaria; el
\cref{thm:shift-hensel-completo} proporciona la única memoria 3-ádica que
emite sus bloques. La continuidad se comprueba por cilindros: fijar \(T\)
bloques fija \(6T\) cifras y encierra el valor en un intervalo de diámetro
\(3^{-6T}\). La identificación del cociente con la imagen es la factorización
canónica de una aplicación sobreyectiva por sus fibras.
\end{proof}

Una sola carta 3-ádica es compacta y, por tanto, no puede cubrir de manera
continua una recta no acotada. La globalización correcta utiliza una familia
numerable de cartas.

\begin{corollary}[Cobertura arquimediana numerable de \(\mathbb R\)]
\label{cor:atlas-hmt-reales}
Sea
\[
 \widehat X_{\rm raw}=\mathbb Z\times\mathbb Z_3^6,
 \qquad
 \widehat P_{\rm raw}(m,x)=m+P_{\rm raw}(x).
\]
Entonces \(\widehat P_{\rm raw}\) es sobreyectiva sobre \(\mathbb R\).
En consecuencia, todo real admite al menos una genealogía Hensel bruta y
\[
 \widehat X_{\rm raw}/\!\sim_{\widehat P_{\rm raw}}
 \ \cong\ \mathbb R
\]
como conjunto de valores.
\end{corollary}

\begin{proof}
Para \(r\in\mathbb R\), tome \(m=\lfloor r\rfloor\) y aplique el
\cref{thm:cobertura-arquimediana-bruta} a \(r-m\in[0,1)\).
\end{proof}

Este resultado precisa la afirmación según la cual \(\mathbb R\) es la imagen
arquimediana cociente de
una estructura genealógica: la proyección olvida la memoria 3-ádica y conserva
el punto evaluado. La afirmación demostrada aquí concierne a la completación
\emph{bruta}. Si \(X_{\rm surv}\subseteq\mathbb Z_3^6\) denota el subespacio
que además satisface una familia concreta de cierres APP--TRIT--TPK, entonces
\[
 P_{\rm raw}(X_{\rm surv})=[0,1]
\]
es un teorema adicional. El transductor garantiza capacidad completa; la
supervivencia determina qué parte de esa capacidad realiza el sistema
restringido.

Tampoco se deduce aquí que la suma y el producto nativos de TPK desciendan a
todo el cociente global. Se puede transportar la aritmética usual desde
\(\mathbb R\), pero demostrar que coincide con operaciones construidas antes
de evaluar exige el correspondiente teorema de descenso.

En esta construcción, los cilindros no son infinitesimales no nulos de
\(\mathbb R\). Cada cilindro de profundidad finita tiene diámetro positivo,
y la sucesión de diámetros converge a cero. Dos secciones que coinciden hasta
la profundidad \(T\) y difieren después forman una \emph{perturbación de
cola}; su diferencia arquimediana queda acotada por el diámetro del cilindro
común. Ésta es la formulación exacta de los ``cilindros que tienden a cero''
dentro de la recta real ordinaria.

\section{Carta posicional en base \(10^3\) y clasificación de las secciones publicadas}

La completación cúbica
\(1000=729+243+27+1\) hace natural agrupar las cifras decimales en bloques
\(b_n\in\{0,\ldots,999\}\). Para una cinta de bloques definimos
\[
 \operatorname{ev}_{1000}(b_1,b_2,\ldots)
 =\sum_{n\geq1}\frac{b_n}{1000^n}.
\]
La expansión canónica excluye una cola finalmente constante de bloques
\(999\). Esta convención resuelve la duplicidad de los extremos de cilindro.

\begin{theorem}[Taxonomía visible en base mil]
\label{thm:taxonomia-base-mil}
Sea \(x\in[0,1)\) y sea \((b_n)\) su expansión canónica en base mil.
Entonces:
\begin{enumerate}
\item \(x\) posee una expansión terminante si y sólo si, escrito como
fracción irreducible \(a/q\), se tiene \(q\mid1000^N\) para algún \(N\);
equivalentemente, los únicos primos que dividen \(q\) son \(2\) y \(5\).
\item \(x\) es racional si y sólo si \((b_n)\) es finalmente periódica.
\item \(x\) es irracional si y sólo si \((b_n)\) no es finalmente
periódica.
\end{enumerate}
\end{theorem}

\begin{proof}
Una expansión que termina en profundidad \(N\) tiene la forma
\(D/1000^N\); al reducirla, su denominador sólo contiene los primos \(2\) y
\(5\). Recíprocamente, si \(q\mid1000^N\), entonces \(a/q\) se escribe con
\(N\) bloques.

Si después de un prefijo de longitud \(r\) se repite un periodo de longitud
\(s\) cuyo entero de base mil es \(P\), la cola es una serie geométrica:
\[
 \frac{P}{1000^r(1000^s-1)},
\]
y el valor es racional. A la inversa, al expandir \(a/q\), los restos de la
división sucesiva sólo pueden tomar \(q\) valores. Un resto nulo produce una
expansión terminante; la repetición de un resto produce un periodo. La tercera
afirmación es la negación exacta de la segunda.
\end{proof}

La palabra \emph{fracción} designa una presentación \(a/q\) de un racional,
no una especie numérica adicional. Tampoco deben confundirse periodicidad y
algebraicidad. Un real es algebraico cuando anula algún polinomio no nulo de
\(\mathbb Q[X]\); es trascendente cuando no lo hace. Esta segunda
clasificación es independiente de la base usada para escribir las cifras.
Por ejemplo,
\[
 \varphi^2-\varphi-1=0
\]
hace de \(\varphi\) un irracional algebraico, mientras que la trascendencia
de \(e\) y \(\pi\) procede de teoremas clásicos externos a HMT\@.

\section{Concatenación, reducción nonádica y acarreo}

La compatibilidad entre la carta cúbica y la raíz digital es una identidad
aritmética, no una semejanza visual. Si \(u_j\in\{0,\ldots,999\}\) y
\[
 D_n=1000D_{n-1}+u_n,
 \qquad D_0=0,
\]
entonces \(D_n\) es el entero obtenido al concatenar los primeros \(n\)
bloques.

\begin{theorem}[Congruencia nonádica de control de la carta cúbica]
\label{thm:checksum-nonadico-base-mil}
Para todo \(n\geq1\),
\[
 D_n\equiv\sum_{j=1}^n u_j\pmod9.
\]
Por tanto, concatenación en base mil, reducción módulo nueve y raíz digital
son compatibles.
\end{theorem}

\begin{proof}
Como \(1000\equiv1\pmod9\), la recurrencia da
\(D_n\equiv D_{n-1}+u_n\pmod9\). La iteración desde \(D_0=0\) produce la
fórmula.
\end{proof}

En el cierre dodecafásico de \(\alpha\), la ecuación de acarreo por bloque
tiene la forma
\[
 P_m+E_m-\Phi_m-K_m+c_{m+1}=A_m+1000c_m.
\]
Reducida módulo nueve,
\[
 A_m\equiv
 P_m+E_m-\Phi_m-K_m+c_{m+1}-c_m
 \pmod9.
\]
Si los acarreos extremos se anulan, la suma telescópica elimina la memoria
intermedia. El residuo global de la salida coincide entonces con el residuo
de la combinación de entradas. Esta ley explica por qué el registro nonádico
controla la concatenación decimal; no selecciona por sí sola los bloques que
deben concatenarse.

\section{Coordenada nonádica elevada y retorno de fase}

La reducción módulo nueve no debe confundirse con la identificación de los
enteros que poseen el mismo residuo. Sea
\[
 I_9=\{1,2,\ldots,9\}.
\]
Para cada entero positivo \(n\), definimos su \emph{fase nonádica} y su
\emph{índice de vuelta} por
\begin{equation}
 \rho_9(n)=1+((n-1)\bmod9),
 \qquad
 \kappa_9(n)=\frac{n-\rho_9(n)}9.
 \label{eq:coordenada-nonadica-levantada}
\end{equation}
La primera coordenada es la raíz digital escrita en el alfabeto
\(I_9\), de modo que la clase residual cero se representa por la fase \(9\).
La segunda registra cuántas vueltas completas han precedido a la fase
visible.

\begin{theorem}[Descomposición nonádica con memoria]
\label{thm:descomposicion-nonadica-memoria}
La aplicación
\[
 \Theta_9:\mathbb N_{>0}\longrightarrow\mathbb N_0\times I_9,
 \qquad
 n\longmapsto\bigl(\kappa_9(n),\rho_9(n)\bigr),
\]
es biyectiva, con inversa
\[
 \Theta_9^{-1}(q,r)=9q+r.
\]
En estas coordenadas, el sucesor aritmético queda conjugado con la dinámica
\begin{equation}
 \mathcal S_9(q,r)=
 \begin{cases}
  (q,r+1),&1\leq r<9,\\
  (q+1,1),&r=9.
 \end{cases}
 \label{eq:sucesor-nonadico-levantado}
\end{equation}
En particular, el paso visible \(9\to1\) es un retorno de fase acompañado
por el incremento irreversible \(q\to q+1\); no es una igualdad entre los
enteros \(9\) y \(10\).
\end{theorem}

\begin{proof}
La división euclídea de \(n-1\) por \(9\) proporciona de manera única
\(n-1=9q+s\), con \(q\in\mathbb N_0\) y \(0\leq s<9\). Tomando
\(r=s+1\) se obtiene \(n=9q+r\), con \(r\in I_9\), lo que demuestra a la vez
existencia, unicidad y la fórmula inversa. Si \(r<9\), sumar una unidad sólo
incrementa \(r\). Si \(r=9\), la igualdad
\(9q+9+1=9(q+1)+1\) da la segunda rama de
\cref{eq:sucesor-nonadico-levantado}.
\end{proof}

Esta elevación separa tres nociones que la cifra aislada mezcla:
\begin{enumerate}
\item la unidad de avance, realizada por el sucesor;
\item la fase visible \(r\), que recorre el ciclo nonádico;
\item la memoria acumulada \(q\), que distingue retornos de una misma fase.
\end{enumerate}
La proyección \((q,r)\mapsto r\) olvida el número de vueltas y satisface
\[
 \rho_9(n+9)=\rho_9(n),
 \qquad
 \kappa_9(n+9)=\kappa_9(n)+1.
\]
Éste es el modelo aritmético elemental del mecanismo que, en el estado TPK
enriquecido, conserva además cocientes de Hensel, orientación, firma y dato
de frontera. Un retorno de palabra visible puede coexistir así con una
genealogía distinta.

El cero no pertenece a \(I_9\): se añade como punto base separado antes de
aplicar \(\Theta_9\). Para los enteros negativos se conserva el módulo
positivo de la descomposición y se adjunta una orientación
\(\varepsilon\in\{-1,+1\}\). En consecuencia, ni el signo ni el cero se
deducen de la fase residual. Esta tipificación evita convertir una carta
cíclica en una identidad aritmética falsa.

\begin{theorem}[Entero orientado y signo genealógico]
\label{thm:entero-orientado-signo-genealogico}
La aplicación
\[
\Theta_9^\pm:
\mathbb Z\setminus\{0\}
\longrightarrow
\{-1,+1\}\times\mathbb N_0\times I_9,
\]
\[
n\longmapsto
\bigl(\operatorname{sgn}(n),
      \kappa_9(|n|),
      \rho_9(|n|)\bigr)
\]
es biyectiva, con inversa
\[
(\varepsilon,q,r)\longmapsto\varepsilon(9q+r).
\]
La negación aritmética actúa por
\[
(\varepsilon,q,r)\longmapsto(-\varepsilon,q,r):
\]
invierte la orientación y conserva magnitud, fase y memoria de vueltas.
\end{theorem}

\begin{proof}
El \cref{thm:descomposicion-nonadica-memoria} proporciona de manera única
\(|n|=9q+r\), con \(q\in\mathbb N_0\) y \(r\in I_9\). El signo de un entero
no nulo es asimismo único. La fórmula inversa recompone \(n\), y cambiar
\(n\) por \(-n\) sólo cambia \(\varepsilon\).
\end{proof}

Así, un negativo no es una fase residual especial ni una cantidad
``imposible''. Es la misma magnitud nonádica transportada con orientación
genealógica opuesta. Esta estructura será prolongada en Mecánica Dimensional
por la anti-sección partícula--antipartícula.

Iterar \(\kappa_9\) produce una torre finita para todo entero ordinario; en
la completación inversa, una sucesión compatible de cocientes puede ser
infinita. La diferencia entre ambos casos anticipa la taxonomía HMT: el valor
visible responde a la expansión posicional, mientras que el tipo genealógico
responde a la evolución completa de sus memorias.

\section{2 ejes de clasificación: valor visible y genealogía de transporte}

La taxonomía precedente clasifica el valor proyectado. HMT añade un segundo
eje que la escritura decimal ordinaria descarta. Para una sección
\(x\in\Omega_\infty^{\rm Arch}\) distinguimos:
\[
 \begin{array}{rcl}
 \text{tipo visible}&=&
 \text{terminante, periódico o no periódico; algebraico o trascendente},\\
 \text{tipo genealógico}&=&
 \text{clase de acarreo, orientación, cierre y monodromía de }x.
 \end{array}
\]
Dos secciones pueden compartir el mismo real y diferir en el segundo eje. Un
real racional posee una expansión visible finalmente periódica, pero una
presentación enriquecida redundante puede conservar una memoria interna no
periódica que la evaluación olvida. La periodicidad del \emph{valor} se
refiere siempre a su expansión canónica; la periodicidad de la
\emph{genealogía} se refiere al estado enriquecido de TPK.

Esta separación proporciona una definición concisa:
\begin{equation}
 \boxed{
 \text{número HMT}
 =
 \text{sección compatible con memoria};
 \qquad
 \text{valor real}
 =
 \text{su clase de evaluación}.}
 \label{eq:numero-hmt-y-valor-real}
\end{equation}
Medir en esta carta es acoplar la sección al lector, fijar un cilindro y
condicionar las prolongaciones compatibles. Aumentar la precisión significa
estrechar el cilindro, no consumir la genealogía que permanece en la fibra.

\section{Memoria de profundidad arbitraria en sistemas de emisión finita}

La recurrencia visible de nueve fases puede coexistir con una salida
irracional sólo si el estado completo conserva información que no se reduce a
esas nueve etiquetas. El siguiente resultado hace exacta esa necesidad.

\begin{theorem}[Caracterización mediante emisores finitos]
\label{thm:emisor-finito-racional}
Sea \(Q\) un conjunto finito, \(F:Q\to Q\) una transición determinista,
\(o:Q\to\{0,\ldots,999\}\) una salida y \(q_0\in Q\). La cinta
\(b_n=o(F^{n-1}(q_0))\) es finalmente periódica y su evaluación es racional.
Recíprocamente, toda cinta finalmente periódica admite un emisor determinista
finito.
\end{theorem}

\begin{proof}
Entre \(|Q|+1\) estados consecutivos aparecen dos iguales. Desde la primera
repetición, el determinismo obliga a repetir todo el futuro, por lo que la
cinta es finalmente periódica. El \cref{thm:taxonomia-base-mil} da la
racionalidad. Para la recíproca basta construir una cadena de estados para el
prefijo y un ciclo para el periodo.
\end{proof}

Por tanto, una ley que genere exactamente \(e\), \(\pi\) o cualquier otro
irracional no puede consistir únicamente en un autómata determinista de
estado finito que se repite sin memoria adicional. Puede usar un límite
inverso, una memoria 3-ádica, una profundidad creciente o una entrada
aperiódica; todos esos mecanismos proporcionan estado efectivo no acotado.
En HMT, el retorno visible
\(g_{\rm vis}(n+9)=g_{\rm vis}(n)\) no fuerza el retorno del acarreo, de la
orientación ni de la frontera. Esa diferencia es precisamente el lugar donde
puede residir una expansión no periódica.

\section{El cuadrado local de \texorpdfstring{\(e\)}{e} y la ruptura de memoria}

Para una palabra decimal \(a_1\cdots a_n\), llamaremos \emph{cuadrado local}
de longitud \(\ell\) en la posición \(i\) a una factorización
\[
 a_i\cdots a_{i+2\ell-1}=vv,
 \qquad |v|=\ell.
\]
Esta noción pertenece a la combinatoria de palabras. No afirma que la cola
infinita posea un período.

\begin{theorem}[Caracterización finita del cuadrado \(1828^2\)]
\label{thm:cuadrado-local-e}
Considérense las \(243\) palabras del catálogo y el lector fijado
\(w_6\mapsto w_{30}\mapsto d_1d_2d_3d_4\), con cuatro tríadas en base mil.
Existe una única palabra cuya lectura contiene un cuadrado de longitud cuatro
iniciado en la segunda cifra decimal. Es
\[
 w_e=201101,
 \qquad
 718281828459=7\,\underbrace{1828\,1828}_{(1828)^2}\,459.
\]
La repetición no se prolonga a una tercera copia: la cifra que exigiría esa
continuación sería \(1\), mientras que la cifra efectiva es \(4\). En todo el
catálogo sólo existe otro cuadrado de longitud cuatro,
\[
 575157518472=\underbrace{5751\,5751}_{\text{cuadrado local}}\,8472,
\]
que comienza en la primera cifra.
\end{theorem}

\begin{proof}
La enumeración exhaustiva evalúa las \(243\) palabras mediante aritmética
entera. El perfil de cuadrados para longitudes \(1,\ldots,6\) contiene
\[
 240,\ 21,\ 3,\ 2,\ 0,\ 0
\]
ocurrencias, distribuidas entre
\[
 153,\ 19,\ 3,\ 2,\ 0,\ 0
\]
palabras. Las dos ocurrencias de longitud cuatro son exactamente las del
enunciado. Dos enumeraciones enteras independientes reproducen el mismo censo
y sus salidas completas coinciden término a término.
\end{proof}

La palabra \(w_e\) posee además la cadena de elevación
\[
201101\mid121221\mid102011\mid012222\mid102011.
\]
El tercer y el quinto bloques visibles coinciden. No coinciden, sin embargo,
los estados que los transportan. Antes de ambas emisiones, sus reducciones
módulo \(27\) son
\[
(25,11,7,17,8,16),
\qquad
(7,8,4,23,26,7),
\]
y los cocientes de Hensel posteriores son
\[
(6,13,9,2,13,1),
\qquad
(15,0,3,19,23,25).
\]
Por tanto, el retorno \(102011\) es un retorno de la proyección visible, no
del estado enriquecido. Ésta es la formulación exacta de la repetición local
que se rompe: el cociente conserva una genealogía que el bloque residual no
registra.

En consecuencia, el comienzo
\[
e=2.718281828459\ldots
\]
no presenta una ``semiperiodicidad'', sino un cuadrado local de palabra con
retorno visible y ruptura de memoria. El resultado es exhaustivo dentro del
lector fijado; su contenido es combinatorio y dinámico, no una inferencia de
irracionalidad a partir de doce cifras.

\section[2 medidas finitas del emisor 468 a 243]{2 medidas finitas del emisor \texorpdfstring{\(468\to243\)}{468→243}: frecuencia uniforme y prioridad estructural}

El catálogo público distingue \(468\) emisiones y \(104\,976\) semillas.
Sea
\[
 q:E_{468}\longrightarrow W_{243}
\]
la aplicación que olvida la firma de emisión y conserva la palabra de seis
trits. Hay dos medidas naturales, porque hay dos universos finitos de conteo:
\[
 \mu_{468}(w)=\frac{|q^{-1}(w)|}{468},
 \qquad
 \mu_{\rm seed}(w)=
 \frac{\sum_{e\in q^{-1}(w)}m(e)}{104\,976},
\]
donde \(m(e)\) es la multiplicidad de semillas de la emisión \(e\).

\begin{theorem}[Medidas canónicas relativas al universo contado]
\label{thm:medidas-emisor-468}
La medida uniforme sobre \(E_{468}\) es la única probabilidad invariante bajo
todas las permutaciones de sus emisiones; su medida imagen es \(\mu_{468}\).
La medida uniforme sobre las \(104\,976\) semillas induce
\(\mu_{\rm seed}\). Sus histogramas exactos son
\[
\begin{array}{c|rrrrrr}
|q^{-1}(w)|&1&2&3&4&5&6\\ \hline
\#\{w\}&132&66&18&3&6&18
\end{array}
\]
y
\[
\begin{array}{c|rrrrrrrrr}
\text{peso de semillas}
&144&288&432&576&864&1008&1440&1944&2088\\ \hline
\#\{w\}
&120&54&18&9&15&6&3&12&6.
\end{array}
\]
Para las tres palabras distinguidas por el lector se obtiene
\[
\begin{array}{c|cc}
&\mu_{468}&\mu_{\rm seed}\\ \hline
\pi&5/468&7/729\\
e&1/468&1/729\\
\varphi&1/468&1/729.
\end{array}
\]
\end{theorem}

\begin{proof}
La invariancia bajo el grupo simétrico hace iguales las masas de todos los
puntos de \(E_{468}\); la normalización las fija en \(1/468\). La fórmula de
la medida imagen resulta al sumar sobre cada fibra. El mismo argumento sobre el
conjunto de semillas produce la segunda medida. Los histogramas se obtienen
por enumeración exhaustiva del catálogo finito de \(468\) emisiones y \(243\)
palabras; sus sumas verifican
\[
132+66+18+3+6+18=243,
\]
\[
132+2\cdot66+3\cdot18+4\cdot3+5\cdot6+6\cdot18=468
\]
y la suma total de los pesos es \(104\,976\).
\end{proof}

La no uniformidad de la medida imagen es un teorema finito: distintas palabras
tienen distintos números de realizaciones. No existe, en cambio, una
probabilidad uniforme normalizable sobre toda \(\mathbb R\), y estas medidas
no ordenan intrínsecamente todos los reales. La palabra de \(\pi\) pertenece
a una de las seis fibras de tamaño cinco y peso \(1008\); no es el máximo
único de ninguno de los dos histogramas. Su singularidad demostrada procede
de la estructura compuesta de sus cinco regiones, de su partición orientada
\(3+2\) y de sus compatibilidades posteriores, no de declarar a priori una
probabilidad absoluta sobre el continuo.

\section{Selector orbital finito anterior al lector arquimediano}
\label{sec:selector-orbital-constantes}

Las medidas anteriores revelan multiplicidades, pero una palabra aislada de
seis trits todavía depende de una orientación. El catálogo posee una simetría
interna que permite seleccionar primero clases no orientadas y fijar después
un representante. Escribamos una palabra como \(abc\mid def\) y definamos
\[
 r(abc\mid def)=cab\mid efd,
 \qquad
 s(abc\mid def)=def\mid abc.
\]
Estas permutaciones actúan también sobre las seis coordenadas del dato
\(U_6\). En el catálogo completo satisfacen
\[
 r^3=1,
 \qquad s^2=1,
 \qquad srs=r^{-1},
\]
por lo que generan una acción que denotamos \(D_3^{\rm cat}\). La acción usa
el etiquetado TRIT fijo
\[
 \Psi(U_6)=(-U_6\bmod3)
\]
y la partición publicada de las seis coordenadas en dos tríadas.

Para una palabra \(w\), sean \(f(w)=|q^{-1}(w)|\) el número de emisiones,
\(m(w)\) su masa de semillas y
\[
 \nu(w)=(n_0(w),n_1(w),n_2(w))
\]
su censo ordenado de trits. Sea además \(d(e)\) la clase diagonal aditiva de
una emisión y \(o(e)\in\{\mathrm{ES},\mathrm{NO}\}\) su orientación
publicada.

\begin{theorem}[Selección orbital interna de las tres palabras]
\label{thm:selector-orbital-constantes}
La acción \(D_3^{\rm cat}\) preserva el conjunto de las \(468\) emisiones,
la proyección \(q:E_{468}\to W_{243}\) y la multiplicidad de semillas. Su
acción sobre \(W_{243}\) tiene exactamente \(43\) órbitas: \(38\) de tamaño
seis y cinco de tamaño tres. Dentro de ese cociente se cumplen las siguientes
unicidades.

\begin{enumerate}
\item Existe una única órbita \(\mathcal O_\pi\) de tamaño seis tal que, para
todo \(w\in\mathcal O_\pi\),
\[
 f(w)=5,
 \qquad m(w)=1008,
 \qquad
 \{m(e):e\in q^{-1}(w)\}=\{432,144,144,144,144\}.
\]
Es
\[
\mathcal O_\pi=
\{001112,010211,100121,112001,121100,211010\},
\]
y todas sus palabras tienen censo \(\nu=(2,3,1)\).

\item Consideremos las órbitas de tamaño seis cuyas palabras tienen una sola
emisión, de multiplicidad \(144\), y cuyo soporte diagonal conjunto es
\(\{0,3,6\}\). Hay exactamente seis. Entre ellas existe una única órbita
con el mismo censo \((2,3,1)\) de \(\mathcal O_\pi\):
\[
\mathcal O_e=
\{011120,012110,101201,110012,120011,201101\}.
\]

\item En el mismo sector existe una única órbita equilibrada, es decir, con
\(\nu=(2,2,2)\):
\[
\mathcal O_\varphi=
\{002112,020211,112002,121200,200121,211020\}.
\]
\end{enumerate}

Finalmente, la cláusula orientadora
\begin{equation}
 \exists e\in q^{-1}(w):
 \qquad d(e)=6,
 \qquad o(e)=\mathrm{ES}
 \label{eq:gauge-orbital-constantes}
\end{equation}
selecciona un único representante de cada órbita y produce, respectivamente,
\[
 \boxed{010211,\qquad201101,\qquad121200.}
\]
\end{theorem}

\begin{proof}
Las relaciones de grupo se verifican directamente como identidades de
permutaciones de seis coordenadas. La enumeración de las \(468\) filas
comprueba que \(r\) y \(s\) envían cada \(U_6\) a otra fila del catálogo con
la misma multiplicidad, y que la proyección ternaria conmuta con ambas
acciones. Por tanto, las órbitas pueden calcularse sobre las \(243\) palabras
sin perder las fibras.

La partición exhaustiva da el histograma de tamaños
\[
 38\cdot6+5\cdot3=243.
\]
Al filtrar esas \(43\) órbitas por el primer predicado queda exactamente una;
al filtrar por fibra unipuntual, multiplicidad mínima y soporte diagonal
\(\{0,3,6\}\) quedan seis. El censo \((2,3,1)\) deja una sola órbita en el
primer caso y el censo equilibrado \((2,2,2)\) deja una sola en el segundo.
Por último, la búsqueda de una elevación con \(d=6\) y orientación
\(\mathrm{ES}\) deja una palabra en cada una de las tres órbitas. Todas las
filtraciones se realizan sobre descriptores estructurales del catálogo antes de consultar el
lector decimal.
\end{proof}

Los predicados del teorema pueden reunirse en un invariante de selección.
Para una órbita \(\mathcal O\) definimos su \emph{perfil orbital}
\[
 \Sigma_{\rm cat}(\mathcal O)=
 \bigl(
 |\mathcal O|,\ f,\ m,\ \mathcal M_{\rm fib},\
 \mathcal D,\ \nu
 \bigr),
\]
donde \(f\) es el tamaño de fibra, \(m\) la masa de semillas,
\(\mathcal M_{\rm fib}\) el multiconjunto de multiplicidades de sus
elevaciones, \(\mathcal D\) el soporte diagonal cuando se emplea y \(\nu\)
el censo de trits. Las tres clases distinguidas ocupan posiciones diferentes
en este espacio de invariantes:
\[
\begin{array}{c|c|c|c|c|c}
\mathcal O&|\mathcal O|&f&m&\mathcal M_{\rm fib}&(\mathcal D,\nu)\\ \hline
\mathcal O_\pi&6&5&1008&\{432,144,144,144,144\}&(-,(2,3,1))\\
\mathcal O_e&6&1&144&\{144\}&(\{0,3,6\},(2,3,1))\\
\mathcal O_\varphi&6&1&144&\{144\}&(\{0,3,6\},(2,2,2)).
\end{array}
\]
La singularidad demostrada no significa que estos números reciban una
probabilidad absoluta sobre \(\mathbb R\). Significa que, dentro del universo
finito producido por el emisor y antes de evaluar cifras, sus clases son
separables por invariantes internos y que la clase de \(\pi\) posee una fibra
regional de tipo distinto. Ésta es la formulación matemática de una
prioridad estructural en HMT\@.

El control negativo es informativo. Si se omite el soporte diagonal
\(\{0,3,6\}\), existen cuatro órbitas unipuntuales equilibradas y
\(\mathcal O_\varphi\) deja de ser única. Antes de fijar el calibre hay seis
orientaciones posibles en cada órbita. Después de fijarlo, la palabra de
\(\pi\) conserva tres elevaciones positivas de \(U_6\), no una sola región.
Así, el teorema selecciona la palabra orientada pero respeta la multiplicidad
geométrica de las cinco regiones.

La no uniformidad se refuerza al pasar a órbitas:
\[
\begin{array}{c|cc}
&\text{emisiones}&\text{semillas}\\ \hline
\mathcal O_\pi&5/78&14/243\\
\mathcal O_e&1/78&2/243\\
\mathcal O_\varphi&1/78&2/243.
\end{array}
\]
Estos cocientes siguen siendo masas del catálogo finito, no probabilidades
asignadas a los números reales evaluados.

El avance causal es preciso. Las tres palabras dejan de ser entradas
nombradas: se recuperan desde la simetría y los invariantes finitos de
APP--TPK, relativamente a la acción \(D_3^{\rm cat}\), al etiquetado TRIT,
al reparto \(3+3\) y al calibre orientador de
\cref{eq:gauge-orbital-constantes}. La prueba no usa \(L_0\), \(L_1\),
cifras decimales, \(\alpha\), CODATA ni el estado dodecafásico. Siguen siendo
afirmaciones posteriores la unicidad absoluta de esa acción entre todos los
automorfismos posibles, la deducción del calibre en vez de su declaración, la
canonicidad independiente del valor objetivo del lector \(L_0/L_1\) y la emisión coinductiva de las
expansiones infinitas.

\section{Posición de \texorpdfstring{\(\pi,e,\varphi\) y \(\alpha\)}{pi, e, phi y alpha} en la clasificación}

El \cref{thm:shift-hensel-completo} explica por qué cualquier prefijo
decimal puede reemitirse desde un estado 3-ádico: la capacidad bruta es
universal. Por tanto, la mera operación de ida y vuelta sobre una cinta no selecciona una
constante. El contenido distintivo de la construcción publicada se encuentra
en los objetos anteriores y posteriores a esa reemisión:
\begin{enumerate}
\item las fibras finitas y las cinco regiones distintas de \(\pi\);
\item el lector arquimediano tipado que asocia palabras, cilindros y bloques;
\item la dodecafase firmada y el cierre reversible
\[
U_{12}^{\rm sgn}
\longleftrightarrow K
\longleftrightarrow(D_3K,D_4K,Q)
\longleftrightarrow\alpha_{\rm HMT};
\]
\item los criterios de contraste que separan una equivalencia de cartas de un
selector anterior a la evaluación.
\end{enumerate}

Las doce ternas decimales publicadas de \(\pi,e,\varphi\) son coordenadas
finitas exactas del lector enriquecido. La cadena impresa de treinta y seis
cifras es la proyección dodecafásica racional
\(\alpha_{12}=\pi_{12}(\alpha_{\rm HMT})\):
\[
\frac{7297352569283800997285105472380663}{10^{36}}.
\]
La identidad exacta en esta ventana es
\[
\alpha_{12}
:=\pi_{12}(\alpha_{\rm HMT})
=\frac{7297352569283800997285105472380663}{10^{36}}.
\]
La recurrencia compatible a profundidad arbitraria construye la sección
infinita
\[
 \alpha_{\rm HMT}:=\varprojlim_N A^{(N)},
\]
y cada ventana racional es una proyección exacta de ese mismo objeto. El
núcleo \(E_{21/22}\) construye, en otra carta, la raíz positiva simple
finita \(\alpha_E\); el transporte graduado del TPK prolonga ese núcleo en
una familia compatible cuya raíz límite es \(\alpha_B\). Si
\(\alpha_A:=\varprojlim_N\rho_N(\alpha_{12})\), los cuadrados de truncamiento
de las dos familias conmutan y determinan en cada profundidad el mismo
cilindro superviviente. El teorema exacto de las dos vías establece, por
tanto,
\[
 \boxed{\alpha_A=\alpha_B=\alpha_{\rm HMT}}.
\]
La igualdad de \(\operatorname{Tr}_9(\alpha_E^{-1})\) y
\(\operatorname{Tr}_9(\alpha_{12}^{-1})\) es el primer certificado finito de
esa identidad; no la define, no limita su alcance y no selecciona las
prolongaciones. La proyección tipada conserva distinguidas la sección
coinductiva, su ventana racional y la raíz analítica sin separar su identidad
en el límite.

La raíz cuadrada de \(-1\) pertenece a otra carta. En el
\cref{chap:algebras-cuadraticas-orden-nueve} se construye
\[
\mathbb F_9=\mathbb F_3[\iota]/(\iota^2+1),
\qquad \iota^2=-1,
\]
y se demuestra que la matriz Witt realiza esa estructura cuadrática. Se trata
de una extensión algebraica finita exacta, no de una categoría visible de
expansiones reales.

\section{Carta de Hensel no filtrada, supervivencia y frontera global}

Los siguientes controles son falsadores de los resultados finitos
precedentes:
\begin{enumerate}
\item que dos estados distintos módulo \(3^T\) emitan la misma palabra de
\(T\) bloques, o que exista una palabra sin preimagen;
\item que falle la inversa explícita de
\cref{eq:inversa-hensel-finita};
\item que aparezca un racional cuya expansión canónica no sea finalmente
periódica, o un irracional cuya expansión sí lo sea;
\item que un emisor determinista finito produzca una cinta no finalmente
periódica;
\item que los censos del catálogo no sumen \(468\), \(243\) y \(104\,976\),
o que las medidas de las tres palabras no coincidan con la enumeración.
\end{enumerate}

La carta Hensel bruta
\[
\mathbb Z_3^6\longrightarrow[0,1]
\]
es continua y sobreyectiva; al añadir la parte entera, su atlas bruto cubre
\(\mathbb R\). Este resultado exacto no implica que toda fibra bruta contenga
un representante que sobreviva a los filtros adicionales. Las ventanas
finitas publicadas y los lectores nombrados de \(\pi,e,\varphi\) se verifican
dentro de HMT, relativamente a sus semillas, regiones,
orientaciones y reglas enriquecidas.

\ifdefined\HMTInternalTraceability
Para un prefijo filtrado, sea \(T_u\) el árbol de prolongaciones compatibles.
Si \(T_u\) es finitamente ramificado y posee un nodo en toda profundidad, el
lema de König produce una rama infinita; si, además, cada nivel estable es
unitario, la rama es única. Este enunciado es condicional: no se ha demostrado
la no vacuidad universal de \(T_u\) para todo prefijo. Permanecen, por tanto,
como \textsc{programa abierto} la cobertura de todas las fibras por
supervivientes filtrados y el descenso bien definido de la suma y el producto
nativos al cociente arquimediano.
\fi
